\section{The Controllers}
\label{sec:controllers}

We implemented 13 controllers (Tbl.~\ref{tab:controllers}), each meeting the platform through the same interface: it reads the current state and the physical limits, and returns a velocity command.
Holding that interface fixed keeps the comparison on the control law.
All numerical parameters, training configurations, and reward functions are detailed in the supplementary material.

\begin{table}[t]
\centering
\caption{The 13 benchmarked controllers. Paradigm labels follow
  standard taxonomy in control and RL literature. The taxonomy has two axes:
  data/model source (analytical model, simulation, hardware logs) and
  decision mechanism (feedback law, MPC/planning, online RL policy, offline RL policy).
  Suffixes: WO = wall override, DR = domain randomization, BZ = blind-zone sensor model,
  WM = learned world model.}
\label{tab:controllers}
\setlength{\tabcolsep}{3pt}
\footnotesize
\begin{tabular}{lll}
\toprule
\textbf{Controller} & \textbf{Data source} & \textbf{Decision mechanism} \\
\midrule
PD + WO        & Analytical model     & Error feedback + tuned gains \\
LQR + WO       & Analytical model     & PD-equivalent state feedback \\
SMC + WO       & Analytical model     & Sliding-mode feedback \\
\midrule
MPC + WO       & Analytical model     & Constrained QP planning \\
NMPC           & Analytical model     & Nonlinear MPC planning \\
\midrule
PPO (BZ+DR)    & Simulation           & Model-free RL (on-policy) \\
TRPO (DR)      & Simulation           & Model-free RL (on-policy) \\
SAC (DR)       & Simulation           & Model-free RL (off-policy) \\
TD3 (DR)       & Simulation           & Model-free RL (off-policy) \\
TQC (DR)       & Simulation           & Model-free RL (off-policy) \\
\midrule
PPO (WM)       & Logged data          & RL trained in learned model \\
MPPI           & Logged data          & Sampling MPC via learned model \\
IQL            & Logged data          & Offline model-free RL \\
\bottomrule
\end{tabular}
\end{table}

\subsection{Classical Control (PD, LQR, SMC)}

The three classical controllers follow the standard workflow: model, tune in simulation, deploy.
PD acts on measured position and velocity errors, LQR on the linearized dynamics, and SMC on a nonlinear sliding surface.
We tuned PD and SMC with CMA-ES~\cite{hansen2016cma}, an evolutionary optimizer, in simulation over stratified initial conditions.
For LQR, the CMA-ES-tuned gain does not transfer to the hardware, so we deploy a PD-equivalent gain instead; the supplement's classical-tuning appendix gives the details.
Each law is designed around the equilibrium and carries no notion of the rail ends from Sec.~\ref{sec:system}.

\noindent\textbf{The wall-override mechanism:} To give the nominal controllers a way to respect the boundary, we add a simple heuristic: when the cart approaches the rail limit, or wall, override whatever the nominal controller commands and drive away from it.
PD, LQR, SMC, and MPC all run with this guard, denoted ``+WO'' in our results; NMPC handles the rail through its own constraint formulation and runs without it.

\subsection{Predictive Control (MPC, NMPC)}

Linear MPC solves a constrained quadratic program (QP) at each timestep using linearized dynamics; NMPC uses the full nonlinear model over a horizon of $N=10$ steps.
Both were tuned with the same CMA-ES framework, in this case over cost weights and horizon parameters.
Near the rail, linear MPC's cart-position constraint can make that QP infeasible, so the deployed MPC~(+WO) drops the constraint and falls back on the same wall override as the classical controllers.

\subsection{Simulation-Trained Reinforcement Learning}

We trained five RL algorithms, PPO~\cite{schulman2017proximal}, TD3, TQC, SAC, and TRPO, in simulation with domain randomization.
The training environment includes the velocity interface of Sec.~\ref{sec:system}, so each policy trains against the servo dynamics the hardware presents; all use small two-layer MLPs, with architectures and training budgets in the supplement.
One variant, PPO (BZ+DR), also models the sensor blind zone of Sec.~\ref{sec:system} during training: with the simulated readings freezing near the arc edges as the real ones do, the policy learns to act cautiously while its inputs may already be stale.
We ablate this blind-zone model on PPO alone, comparing PPO with the sensor model against PPO without it (ablations in the supplementary material).

\subsection{Hardware-Data-Driven Controllers (IQL, MPPI, PPO-WM)}

These three controllers all learn from real hardware data, and each puts that data to work in a different way.

\noindent\textbf{Offline RL (IQL):} IQL~\cite{kostrikov2022iql} learns a policy directly from logged hardware data, requiring no simulator.
We trained it on approximately 333,000 transitions (1,415 trials) collected from every controller we ran on hardware (MPPI, NMPC, RL, MPC, PD, LQR, SMC, and earlier offline-RL checkpoints, though not the deployed IQL policy itself).
The dataset is mixed, 83\% successful episodes and 17\% failures, and rewards are recomputed offline from the logged states.
IQL extracts its policy conservatively, through expectile regression~\cite{kostrikov2022iql}, which keeps it clear of the distributional-shift instabilities common to other offline RL methods.

\noindent\textbf{Model-based planning (MPPI):} This approach learns the system's dynamics directly from hardware data: a two-layer LSTM trained to predict state transitions from a 1M-transition random-exploration dataset.
At runtime MPPI treats this learned model as an internal simulator.
Every control step (50\,ms) it samples hundreds of candidate action sequences, rolls each forward through the network, and executes the first action of the cost-weighted-average sequence, replanning continuously as new data arrives.

\noindent\textbf{RL in a learned world (PPO-WM):} A second variant trains a standard PPO policy inside a learned world model~\cite{hafner2020dream} rather than using it for online planning.
The result is a lightweight network that trains and runs on CPU, with no GPU at deployment, while still drawing on dynamics learned from real data.
MPPI and PPO-WM share the same world-model architecture but were trained on two separate 1M-transition collections; the collection behind the PPO-WM model suits RL training better, see reasons in the supplement's world-model appendix.

All thirteen then meet the same evaluation protocol, defined in the next section.
